\documentclass[border=3pt]{standalone}

\usepackage{tikz}
\usepackage{tikz-3dplot}
\usetikzlibrary{arrows.meta,calc,decorations.markings}

\usepackage{physics}
\usepackage{bm}

% Styles
\tikzset{axis/.style={thick,-{Latex[length=3.2mm]}}}
\tikzset{vec/.style={very thick,-{Latex[length=3.2mm]}}}
\tikzset{svec/.style={vec,black}}
\tikzset{bvec/.style={vec,black}}
\tikzset{sdot/.style={vec,red!75!black}}

\begin{document}

\tdplotsetmaincoords{70}{110}

% Angles for the point on the sphere
\pgfmathsetmacro{\thetavec}{48.17} % polar angle (from +z)
\pgfmathsetmacro{\phivec}{63.5}    % azimuth (from +x in xy)

\begin{tikzpicture}[tdplot_main_coords, line cap=round, line join=round, font=\Large]

% Parameters
\def\R{5}
\def\OmegaLen{1.4}

% Origin
\coordinate (O) at (0,0,0);

% Axes (optional)
%\draw[axis] (O) -- (6.5,0,0) node[pos=1.05] {$x$};
%\draw[axis] (O) -- (0,6,0) node[pos=1.05] {$y$};
%\draw[axis] (O) -- (0,0,6) node[pos=1.05] {$z$};

% --- 3D sphere outline (projected circle)
\begin{scope}[tdplot_screen_coords]
  \draw[line width=1.2pt] (0,0) circle (\R);
\end{scope}

% --- Single 3D ellipse: great circle in the xy-plane (equator)
\draw[line width=1.2pt] (0,0,0) circle (\R);

% --- Rotation axis (z) and omega
\draw[dashed, line width=1.1pt] (0,0,-\R) -- (0,0,\R);
\draw[vec] (0,0,\R) -- ++(0,0,\OmegaLen) node[above] {$\omega_{\oplus}$};



% --- Point P on sphere giving s(t) (first quadrant)
\tdplotsetcoord{P}{\R}{\thetavec}{\phivec}
\pgfmathsetmacro{\zP}{\R*cos(\thetavec)}
\pgfmathsetmacro{\rSmall}{\R*sin(\thetavec)}

% Small circle (parallel to xy) passing through P
\draw[line width=1.1pt] (0,0,\zP) circle (\rSmall);

\fill (P) circle (1.2pt);

% --- Vector s(t)
\draw[svec] (O) -- (P) node[midway, below right] {$\mathbf{s}(t)$};

% --- Declination $\delta$ (measured from the xy-plane upward)
\pgfmathsetmacro{\deltavec}{90-\thetavec}
\tdplotsetcoord{Pxy}{\R}{90}{\phivec}
\draw[dashed] (O) -- (Pxy);
\draw[dashed] (P) -- (Pxy);

\tdplotsetthetaplanecoords{\phivec}
\tdplotdrawarc[tdplot_rotated_coords]{(O)}{1.2}{90}{48}{anchor=east}{$\delta$}

% --- b: midpoint at the intersection of the meridian circle with the xy-plane
% (choose the intersection in the direction of azimuth \phivec)
\pgfmathsetmacro{\bx}{\R*cos(\phivec)}
\pgfmathsetmacro{\by}{\R*sin(\phivec)}
\coordinate (bMid) at (\bx,\by,0);

% tangent direction in the xy-plane (same as \dot{s}(t))
\pgfmathsetmacro{\tx}{-sin(\phivec)}
\pgfmathsetmacro{\ty}{cos(\phivec)}

\coordinate (bTail) at ($(bMid)+(-1.8*\tx,-1.8*\ty,0)$);
\coordinate (bHead) at ($(bMid)+(1.8*\tx,1.8*\ty,0)$);
\draw[bvec] (bTail) -- (bHead) node[midway, below right] {$\mathbf{b}$};

% --- \dot{s}(t): tangent to the motion of s(t) (parallel to b), starting at P
\coordinate (sdotHead) at ($(P)+(2.6*\tx,2.6*\ty,0)$);
\draw[sdot] (P) -- (sdotHead)
  node[near end, above right, black] {$\dot{\mathbf{s}}(t)=\omega_{\oplus}\times\mathbf{s}(t)$};

\end{tikzpicture}
\end{document}